\section{Introduction}\label{sec:introduction}

Exact inference in probabilistic graphical models is usually developed for discrete variables. Variable
elimination and belief propagation pass tables of numbers between factors, and summing out a variable is a sum
over its values \cite{koller2009probabilistic}. Many estimation problems in robotics do not fit this setting. The
pose of a robot and the positions of landmarks are continuous, the motion and sensor models are nonlinear, and the
belief about the pose can have several modes, for example when the robot does not yet know where it started.
Summation becomes integration, and for these models the integrals have no closed form.

The particle filter handles this case by representing each belief with a set of weighted samples
\cite{gordon1993novel,arulampalam2002tutorial}. It is approximate inference on a dynamic Bayesian network: the
exact filtering recursion is kept, and each integral is replaced by importance sampling. The approximation is only
as good as the sampling. A filter proposes new particles, weights them, and resamples them, and each of these
steps has alternatives that trade accuracy against particle diversity and runtime. Which alternatives matter,
and by how much, is the question of this report.

This report is organized around the three questions a graphical model raises, and its aim is to answer all three
for one continuous problem. \emph{Representation}: localization and SLAM are written as a dynamic Bayesian
network, and the conditional independence its structure implies, the landmarks being independent of one another
given the trajectory, is exactly what makes Rao-Blackwellized SLAM tractable (\cref{sec:background}).
\emph{Inference}: the particle filter is derived as approximate inference on that network, and the sampling
choices it leaves open are compared experimentally (\cref{sec:particle-filters,sec:experiments}).
\emph{Learning}: the parameters of the network's two conditional distributions, the motion and the measurement
noise, are estimated from a recorded run by maximum likelihood with the trajectory unobserved
(\cref{sec:learning,sec:e11}). Taking the three together is what the graphical-model view contributes here: the
structure decides what can be factorized, the inference decides what can be computed, and the quality of the
inference decides what can be learned.

The study uses two problems. In localization the map is known and the filter estimates the robot's pose from a
uniform start. In simultaneous localization and mapping (SLAM) the filter estimates the pose and a map of
landmarks together; FastSLAM \cite{montemerlo2002fastslam} does this with Rao-Blackwellization, sampling the
robot's trajectory and estimating each landmark analytically given that trajectory. On these two problems the
report compares:
\begin{itemize}
  \item the resampling scheme: multinomial, residual, stratified and systematic resampling;
  \item when to resample: at every step, never, or when the effective sample size or the largest weight shows
        that the weights have degenerated;
  \item the proposal distribution: the motion model against proposals that use the current measurement (the
        auxiliary particle filter, the extended Kalman particle filter, and FastSLAM 2.0);
  \item the number of particles.
\end{itemize}
A last question turns the setting around. Every filter above is told how noisy its motion and its sensors are.
Those numbers are parameters of the two conditional distributions of the network, and the same weights that drive
the filter also estimate how probable the recorded measurements are under them, so they can be learned from a run
instead of being set by hand. The report asks which of them a run actually determines, and what has to be true of
the inference for any of them to be learned.

The filters are built on existing open-source implementations: the particle filter tutorial of Elfring et al.\
\cite{elfring2021particle} and the FastSLAM code of PythonRobotics \cite{sakai2018pythonrobotics}. Their sampling
steps were made interchangeable, and the reused code was checked function by function against the original. That check found errors that would have biased the comparisons, most importantly a FastSLAM
2.0 implementation that never samples from its proposal; the corrections are listed in \cref{app:upstream}.

\Cref{sec:background} writes localization and SLAM as dynamic Bayesian networks, explains Rao-Blackwellization, and
states what learning the noise parameters of such a network means.
\Cref{sec:particle-filters} derives the particle filter and the sampling choices compared here.
\Cref{sec:method} describes the simulated worlds, the implementation and the metrics, \cref{sec:experiments}
presents the results, \cref{sec:discussion} discusses them, and \cref{sec:conclusion} concludes.
